% =====================================================
% 题型：五面体面面垂直与线面角体积综合 (q_tjw_20260606_16_pentahedron.tex)
% =====================================================
% =====================================================
% 难度：★★★ (难题)
% =====================================================
\newcommand{\qSolidPentahedronPropertiesStem}{%
  如图，在五面体 \(ABCDE\) 中，\(EA \perp\) 平面 \(ABC\)，\(CD \parallel AE\)，点 \(D, E\) 在平面 \(ABC\) 的同侧，\(AC \perp BC\)，\(AE = AC = BC = 2CD = 4\)，点 \(M\) 为 \(BE\) 中点。\\
  (1) 求证：平面 \(BDE \perp\) 平面 \(ABE\)；\\
  (2) 求直线 \(ED\) 与平面 \(ABE\) 所成角的余弦值；\\
  (3) 求三棱锥 \(M - ABD\) 的体积。%
}

\newcommand{\qSolidPentahedronPropertiesFigure}[1][1.5]{%
  \begin{tikzpicture}[scale=#1, >=stealth, line join=round, line cap=round]
    % Coordinates (perspective projection)
    \coordinate (A) at (0,0);
    \coordinate (B) at (3,-0.5);
    \coordinate (C) at (1.5,0.8);
    \coordinate (E) at (0,2.5);
    \coordinate (D) at (1.5,2.05); % CD is parallel to AE and half the length
    \coordinate (M) at ($(B)!0.5!(E)$);

    % Draw hidden lines
    \draw[dashed, thick, gray!80] (A) -- (C);
    \draw[dashed, thick, gray!80] (B) -- (C);
    \draw[dashed, thick, gray!80] (C) -- (D);
    \draw[dashed, thick, gray!80] (A) -- (D); % for triangle ABD

    % Draw visible outline and edges
    \draw[thick] (A) -- (B) -- (E) -- cycle;
    \draw[thick] (B) -- (D) -- (E);
    
    % Draw point M and line AM, DM
    \fill (M) circle (1.5pt) node[right, inner sep=2pt] {\small $M$};
    \draw[thick] (A) -- (M);
    \draw[thick] (D) -- (M);

    % Labels
    \node[below left] at (A) {\small $A$};
    \node[below] at (B) {\small $B$};
    \node[right] at (C) {\small $C$};
    \node[above] at (D) {\small $D$};
    \node[above left] at (E) {\small $E$};
    
    \ifteacher
      % 教师端辅助线：建系坐标轴
      \draw[->, dashed, gray, thick] (C) -- ($(C)!1.3!(A)$) node[above, black] {\tiny $x$};
      \draw[->, dashed, gray, thick] (C) -- ($(C)!1.3!(B)$) node[right, black] {\tiny $y$};
      \draw[->, dashed, gray, thick] (C) -- ($(C)!1.3!(D)$) node[left, black] {\tiny $z$};
    \fi
  \end{tikzpicture}%
}

\newcommand{\qSolidPentahedronPropertiesAnswer}{%
  (1) 见解析；(2) \(\dfrac{\sqrt{15}}{5}\)；(3) \(\dfrac{16}{3}\)。%
}

\newcommand{\qSolidPentahedronPropertiesSolution}{%
  (1) 证明：因为 \(EA \perp\) 平面 \(ABC\)，所以 \(EA \perp AC\)，\(EA \perp BC\)。\\
  因为 \(AC \perp BC\)，以 \(C\) 为坐标原点，\(CA, CB\) 所在直线分别为 \(x\) 轴，\(y\) 轴建立空间直角坐标系。\\
  因为 \(EA \perp\) 平面 \(ABC\)，\(CD \parallel AE\)，所以 \(CD \perp\) 平面 \(ABC\)，\(CD\) 所在直线即为 \(z\) 轴。\\
  已知 \(AE = AC = BC = 2CD = 4\)，得各点坐标如下：\\
  \(C(0,0,0)\)，\(A(4,0,0)\)，\(B(0,4,0)\)，\(E(4,0,4)\)，\(D(0,0,2)\)。\\
  点 \(M\) 为 \(BE\) 中点，\(M(2,2,2)\)。\\
  平面 \(ABE\) 内有向量 \(\vec{AB} = (-4,4,0)\)，\(\vec{AE} = (0,0,4)\)。\\
  设平面 \(ABE\) 的法向量 \(\vec{n_1} = (x_1, y_1, z_1)\)，则 \(\begin{cases} -4x_1 + 4y_1 = 0 \\ 4z_1 = 0 \end{cases}\)，取 \(x_1 = 1\)，得 \(\vec{n_1} = (1,1,0)\)。\\
  平面 \(BDE\) 内，\(\vec{DB} = (0,4,-2)\)，\(\vec{DE} = (4,0,2)\)。\\
  设平面 \(BDE\) 的法向量 \(\vec{n_2} = (x_2, y_2, z_2)\)，则 \(\begin{cases} 4y_2 - 2z_2 = 0 \\ 4x_2 + 2z_2 = 0 \end{cases}\)，取 \(y_2 = 1\)，则 \(z_2 = 2, x_2 = -1\)，得 \(\vec{n_2} = (-1,1,2)\)。\\
  计算 \(\vec{n_1} \cdot \vec{n_2} = (1)(-1) + (1)(1) + (0)(2) = 0\)。\\
  所以 \(\vec{n_1} \perp \vec{n_2}\)，故平面 \(BDE \perp\) 平面 \(ABE\)。\\
  【更优解法（利用中点 $M$ 证明线面垂直）：】\\
  由上述建系可知，\(M(2,2,2)\)，\(D(0,0,2)\)，得 \(\vec{DM} = (2,2,0)\)。\\
  又平面 \(ABE\) 内有向量 \(\vec{AB} = (-4,4,0)\)，\(\vec{AE} = (0,0,4)\)。\\
  计算数量积：\\
  \(\vec{DM} \cdot \vec{AB} = 2 \times (-4) + 2 \times 4 + 0 \times 0 = 0 \implies DM \perp AB\)；\\
  \(\vec{DM} \cdot \vec{AE} = 2 \times 0 + 2 \times 0 + 0 \times 4 = 0 \implies DM \perp AE\)。\\
  因为 \(AB \cap AE = A\)，且 \(AB, AE \subset\) 平面 \(ABE\)，\\
  所以 \(DM \perp\) 平面 \(ABE\)。\\
  又因为 \(DM \subset\) 平面 \(BDE\)，\\
  所以平面 \(BDE \perp\) 平面 \(ABE\)。\\
  
  (2) 解：由(1)知，平面 \(ABE\) 的法向量 \(\vec{n_1} = (1,1,0)\)。\\
  直线 \(ED\) 的方向向量 \(\vec{ED} = -\vec{DE} = (-4,0,-2)\)。\\
  设直线 \(ED\) 与平面 \(ABE\) 所成角为 \(\theta\)，则：\\
  \[
  \begin{aligned}
  \sin\theta
  &= \dfrac{|\vec{ED} \cdot \vec{n_1}|}{|\vec{ED}| |\vec{n_1}|} \\
  &= \dfrac{|-4 \times 1 + 0 \times 1 + (-2) \times 0|}
  {\sqrt{(-4)^2+0^2+(-2)^2}\sqrt{1^2+1^2+0^2}} \\
  &= \dfrac{4}{\sqrt{20}\sqrt{2}}
  = \dfrac{4}{2\sqrt{10}}
  = \dfrac{\sqrt{10}}{5}.
  \end{aligned}
  \]
  由于 \(\theta \in \left[0, \dfrac{\pi}{2}\right]\)，所以 \(\cos\theta = \sqrt{1 - \sin^2\theta} = \sqrt{1 - \dfrac{10}{25}} = \dfrac{\sqrt{15}}{5}\)。\\
  即直线 \(ED\) 与平面 \(ABE\) 所成角的余弦值为 \(\dfrac{\sqrt{15}}{5}\)。\\
  
  (3) 解：三棱锥 \(M - ABD\) 的体积 \(V_{M-ABD} = V_{D-ABM}\)。\\
  因为点 \(M\) 是 \(BE\) 的中点，所以 \(\triangle ABM\) 的面积是 \(\triangle ABE\) 面积的一半。\\
  三棱锥 \(M-ABD\) 的体积可通过向量混合积公式求解：\\
  \(\vec{AB} = (-4,4,0)\)，\(\vec{AD} = (-4,0,2)\)，\(\vec{AM} = \vec{OM} - \vec{OA} = (2,2,2) - (4,0,0) = (-2,2,2)\)。\\
  三棱锥 \(M-ABD\) 的体积 \(V = \dfrac{1}{6} \left| (\vec{AB} \times \vec{AD}) \cdot \vec{AM} \right|\)。\\
  \(\vec{AB} \times \vec{AD} = (4 \times 2 - 0 \times 0, 0 \times (-4) - (-4) \times 2, -4 \times 0 - 4 \times (-4)) = (8, 8, 16)\)。\\
  所以 \(V = \dfrac{1}{6} |8 \times (-2) + 8 \times 2 + 16 \times 2| = \dfrac{32}{6} = \dfrac{16}{3}\)。\\
  【等价做法：\(V_{E-ABD} = V_{B-ADE} = \dfrac{1}{3} S_{\triangle ADE} \cdot d_{B \to ADE}\)。\\
  因为 \(CD \parallel AE\)，\(AE \perp\) 底面，四边形 \(ACDE\) 是直角梯形，在 \(xz\) 平面上，故 \(B(0,4,0)\) 到平面 \(ACDE\) 的距离为 \(4\)。\\
  \(S_{\triangle ADE} = \dfrac{1}{2} AE \cdot AC = \dfrac{1}{2} \times 4 \times 4 = 8\)。\\
  故 \(V_{E-ABD} = \dfrac{1}{3} \times 8 \times 4 = \dfrac{32}{3}\)。\\
  因为 \(M\) 为 \(BE\) 中点，\(M\) 到平面 \(ABD\) 距离为 \(E\) 到平面 \(ABD\) 距离的一半，所以 \(V_{M-ABD} = \dfrac{1}{2} V_{E-ABD} = \dfrac{16}{3}\)。】%
}

\newcommand{\qSolidPentahedronPropertiesDiagnosis}{%
  （留白待收集学生作答后补充）%
}

\newcommand{\qSolidPentahedronPropertiesTeachingNotes}{%
  精讲。本题考查建系法解决面面垂直与线面角，以及灵活运用等体积法或坐标混合积求体积。%
}


% =====================================================
% 别名兼容：旧课次宏定义
% =====================================================
\newcommand{\qTJWZeroSixSixteenStem}{\qSolidPentahedronPropertiesStem}
\newcommand{\qTJWZeroSixSixteenFigure}[1][1.2]{\qSolidPentahedronPropertiesFigure[#1]}
\newcommand{\qTJWZeroSixSixteenAnswer}{\qSolidPentahedronPropertiesAnswer}
\newcommand{\qTJWZeroSixSixteenSolution}{\qSolidPentahedronPropertiesSolution}
\newcommand{\qTJWZeroSixSixteenDiagnosis}{\qSolidPentahedronPropertiesDiagnosis}
\newcommand{\qTJWZeroSixSixteenTeachingNotes}{\qSolidPentahedronPropertiesTeachingNotes}
